\chapter{Gestão e Governança de TI}

\section{Governança não é gestão}

\begin{teoria}[distinção central]
\termo{Governança} avalia necessidades e opções das partes interessadas, direciona prioridades/decisões e monitora resultados, conformidade e desempenho. \termo{Gestão} planeja, constrói/adquire, executa e monitora atividades em alinhamento à direção estabelecida.

A diferença aparece em COBIT e na ISO/IEC 38500. Em prova, ``governar'' não significa administrar tickets, configurar servidores ou executar projetos; significa estabelecer direção e accountability para que a tecnologia gere valor dentro de risco e recursos aceitáveis.
\end{teoria}

\section{COBIT 2019}

COBIT 2019 é framework para governança e gestão da informação e tecnologia empresarial. O modelo central contém objetivos de governança e gestão distribuídos em cinco domínios.

\begin{table}[ht]
\centering
\begin{tabularx}{\textwidth}{l l X}
\toprule
Domínio & Natureza & Foco \\
\midrule
EDM & Governança & Evaluate, Direct and Monitor. \\
APO & Gestão & Align, Plan and Organize. \\
BAI & Gestão & Build, Acquire and Implement. \\
DSS & Gestão & Deliver, Service and Support. \\
MEA & Gestão & Monitor, Evaluate and Assess. \\
\bottomrule
\end{tabularx}
\end{table}

\subsection{Sistema de governança}

Componentes incluem processos; estruturas organizacionais; princípios, políticas e procedimentos; informação; cultura, ética e comportamento; pessoas, habilidades e competências; serviços, infraestrutura e aplicações. O sistema deve ser \termo{customizado}: estratégia, objetivos, perfil de risco, ameaças, papel da TI, modelo de sourcing e outros fatores influenciam desenho.

\subsection{Cascata de objetivos}

Necessidades das partes interessadas são traduzidas em objetivos empresariais, objetivos de alinhamento e objetivos de governança/gestão. Essa cascata evita implementar controles porque ``o framework manda'' sem conexão com valor e risco institucional.

\subsection{Capacidade}

COBIT 2019 trabalha com níveis de capacidade de processo de 0 a 5. Capacidade avalia quão bem um processo é implementado/desempenhado; não confunda com maturidade organizacional global.

\begin{cebraspe}
EDM é governança. APO, BAI, DSS e MEA são gestão. ``Monitorar'' aparece tanto em governança quanto em gestão, mas com objetos e responsabilidades distintos; não use uma palavra isolada para classificar uma atividade.
\end{cebraspe}

\section{ISO/IEC 38500}

A ISO/IEC 38500 orienta corpos dirigentes sobre uso eficaz, eficiente e aceitável de TI. A lógica clássica é \termo{avaliar, dirigir e monitorar}. A edição vigente é a ISO/IEC 38500:2024; o edital não fixa ano, portanto a revisão final deve observar a referência atual sem perder o núcleo conceitual.

\section{ITIL 4}

O edital exige expressamente \textbf{ITIL v4}. Embora a família ITIL tenha evoluído e exista uma versão mais nova no mercado em 2026, para esta prova o objeto indicado é ITIL 4, salvo retificação.

\subsection{Sistema de Valor de Serviço}

O Service Value System (SVS) integra:
\begin{itemize}
\item princípios orientadores;
\item governança;
\item cadeia de valor de serviço;
\item práticas;
\item melhoria contínua.
\end{itemize}

As quatro dimensões são: organizações e pessoas; informação e tecnologia; parceiros e fornecedores; fluxos de valor e processos.

\subsection{Princípios orientadores}

\begin{enumerate}
\item Focus on value.
\item Start where you are.
\item Progress iteratively with feedback.
\item Collaborate and promote visibility.
\item Think and work holistically.
\item Keep it simple and practical.
\item Optimize and automate.
\end{enumerate}

\subsection{Cadeia de valor}

Atividades: Plan, Improve, Engage, Design \& Transition, Obtain/Build, Deliver \& Support. Fluxos de valor combinam atividades e práticas para transformar demanda/oportunidade em valor.

\subsection{Práticas que mais geram confusão}

\termo{Incident management} busca restaurar operação normal rapidamente e minimizar impacto. \termo{Problem management} reduz probabilidade/impacto de incidentes ao tratar causas reais/potenciais e erros conhecidos. \termo{Service request management} trata solicitações predefinidas iniciadas por usuários. \termo{Change enablement} maximiza mudanças bem-sucedidas por avaliação de risco, autorização e agenda apropriada — não significa burocratizar toda mudança.

\section{Gestão de riscos: ISO 31000 e COSO}

ISO 31000 trata princípios, estrutura e processo de gestão de riscos. O processo inclui comunicação/consulta, escopo/contexto/critérios, avaliação de riscos (identificação, análise e avaliação), tratamento, monitoramento/revisão e registro/reporte.

COSO ERM integra risco à estratégia e ao desempenho. Em auditoria, memorize menos slogans e entenda a lógica: apetite a risco, objetivos, identificação/avaliação, respostas, informação, monitoramento e cultura/governança.

\section{Gestão de projetos — PMBOK 8ª edição}

O edital especifica \textbf{PMBOK 8ª edição}. Essa edição foi publicada pelo PMI em 2025 e combina princípios, domínios de desempenho e orientação de processos de forma não prescritiva.

Pontos de estudo:
\begin{itemize}
\item foco em valor e resultados;
\item adaptação (tailoring) ao contexto;
\item governança e accountability;
\item integração entre abordagens preditivas, adaptativas e híbridas;
\item planejamento, execução, medição, risco, stakeholders, recursos e entrega;
\item interfaces com IA, PMO e aquisições, conforme escopo atualizado.
\end{itemize}

\begin{atencao}
Não use material antigo que trate o PMBOK como lista fixa de 49 processos da 6ª edição para responder pergunta explicitamente sobre a 8ª. Processos continuam relevantes, mas a estrutura e ênfase mudaram.
\end{atencao}

\section{Scrum e Kanban na gestão}

Na governança, métodos ágeis devem preservar transparência, medição, risco, responsabilidade e alinhamento. ``Ágil'' não significa ausência de documentação ou controle. Backlogs, critérios, DoD, métricas de fluxo e evidências automatizadas podem aumentar auditabilidade.

\section{PETI e PDTI}

PETI orienta estratégia de TI em alinhamento à estratégia institucional. PDTI materializa planejamento de iniciativas, necessidades, capacidades, recursos, contratações e priorização em horizonte definido. Nomenclaturas e modelos variam por organização; a lógica cobrada costuma relacionar \termo{estratégia} a \termo{plano executável}.

Um bom planejamento estabelece objetivos, indicadores, metas, responsáveis, iniciativas, riscos, capacidade e governança de acompanhamento.

\section{KPIs e BSC}

KPI mede desempenho de aspecto crítico. Indicador deve ter definição, fórmula, fonte, periodicidade, responsável e interpretação. BSC organiza objetivos e relações causais em perspectivas; no setor público, adaptações são comuns e missão/valor público podem ocupar posição de destaque.

\begin{exemplo}[indicador ruim]
``Número de incidentes fechados'' isoladamente pode incentivar fechamento precoce. Combine volume com reincidência, tempo de restauração, satisfação, backlog envelhecido e taxa de reabertura.
\end{exemplo}

\section{BPMN e gestão de processos}

BPM busca compreender, modelar, executar, medir e melhorar processos. BPMN usa eventos, atividades, gateways, fluxos de sequência, fluxos de mensagem, pools/lanes e artefatos.

Gateway exclusivo escolhe um caminho conforme condição; paralelo divide/sincroniza fluxos concorrentes; inclusivo permite um ou mais caminhos. Fluxo de sequência não deve atravessar pools; comunicação entre participantes usa fluxo de mensagem.

\section{Compliance e conformidade}

Compliance é aderência a obrigações legais, regulatórias, contratuais e internas. Governança estabelece mecanismos para identificar requisitos, atribuir responsabilidade, evidenciar conformidade e corrigir desvios. ``Estar em conformidade'' não garante automaticamente segurança ou bom desempenho; é um piso, não necessariamente o ótimo.

\section{Cibersegurança e continuidade na governança}

A governança deve definir apetite a risco, responsabilidades, métricas, priorização de controles e resposta. ISO/IEC 27001/27002 apoiam SGSI e controles; ISO 22301 apoia continuidade; NIST fornece estruturas e catálogos. Não trate framework como substituto de julgamento de risco.

\section{LAI e LGPD em TI}

LAI promove transparência e acesso, com hipóteses legítimas de restrição. LGPD protege dados pessoais. Não há conflito absoluto: divulgação deve respeitar base legal, finalidade, necessidade, transparência e restrições aplicáveis. Em projeto de dados públicos, classifique o dado antes de simplesmente publicá-lo.

\section{DAMA-DMBOK e governança de dados}

Governança de dados define autoridade, papéis, políticas, qualidade, metadados e decisões sobre dados. Áreas do DMBOK cobrem governança, arquitetura, modelagem, armazenamento, segurança, integração/interoperabilidade, documentos/conteúdo, dados mestres, DW/BI, metadados e qualidade, entre outras.

\section{Mapa mental}

\mapamental{Governança de TI}{
 child {node[concept]{COBIT} child {node[concept]{EDM}} child {node[concept]{APO/BAI/DSS/MEA}}}
 child {node[concept]{ISO 38500} child {node[concept]{avaliar}} child {node[concept]{dirigir/monitorar}}}
 child {node[concept]{ITIL 4} child {node[concept]{SVS}} child {node[concept]{práticas}}}
 child {node[concept]{Projetos} child {node[concept]{PMBOK 8}} child {node[concept]{Agile}}}
 child {node[concept]{Estratégia} child {node[concept]{PETI/PDTI}} child {node[concept]{KPI/BSC}}}
 child {node[concept]{Dados/Risco} child {node[concept]{DAMA}} child {node[concept]{ISO/COSO}}}
}

\section{Questões por nível}

\begin{questao}{\nivel{N1} 14.1}
No COBIT 2019, qual domínio reúne objetivos de governança?
\begin{enumerate}[label=\Alph*)]
\item APO
\item BAI
\item DSS
\item EDM
\item MEA
\end{enumerate}
\end{questao}
\begin{solucao}{14.1}
\textbf{Gabarito: D.} EDM = Evaluate, Direct and Monitor e concentra objetivos de governança. Os demais domínios são de gestão.
\end{solucao}

\begin{questao}{\nivel{N2} 14.2}
Uma equipe resolve rapidamente incidentes recorrentes sem investigar a causa. Em ITIL 4, qual prática deve atuar para reduzir reincidência e impacto por tratamento de causas e erros conhecidos?
\begin{enumerate}[label=\Alph*)]
\item Incident management apenas.
\item Problem management.
\item Service request management.
\item Relationship management apenas.
\item Deployment management.
\end{enumerate}
\end{questao}
\begin{solucao}{14.2}
\textbf{Gabarito: B.} Incident management restaura serviço; problem management trata causas e recorrência. C atende solicitações; D gerencia relações; E move componentes para ambientes.
\end{solucao}

\begin{questao}{\nivel{N3} 14.3}
Qual situação representa governança, e não gestão operacional?
\begin{enumerate}[label=\Alph*)]
\item restaurar serviço após incidente.
\item instalar patch de segurança.
\item avaliar opções estratégicas de sourcing, dirigir a escolha e monitorar geração de valor e risco.
\item escrever script de backup.
\item cadastrar usuário no diretório.
\end{enumerate}
\end{questao}
\begin{solucao}{14.3}
\textbf{Gabarito: C.} Avaliar, dirigir e monitorar escolhas estratégicas é governança. As demais são atividades de gestão/operação.
\end{solucao}

\begin{questao}{\nivel{N4} 14.4}
Um TCE mede sucesso de TI apenas por percentual de projetos ``no prazo'', embora sistemas entregues tenham baixa adoção e alto custo operacional. Qual melhoria é mais alinhada à boa governança?
\begin{enumerate}[label=\Alph*)]
\item manter só indicador de prazo para evitar subjetividade.
\item criar conjunto balanceado de métricas ligado a objetivos e valor público, incluindo benefícios, adoção, qualidade, risco, custo e desempenho operacional.
\item medir somente quantidade de servidores de TI.
\item substituir todos os projetos por Scrum.
\item eliminar metas para impedir gaming.
\end{enumerate}
\end{questao}
\begin{solucao}{14.4}
\textbf{Gabarito: B.} Governança deve medir valor e resultados, não apenas entrega temporal. A cria visão míope; C é métrica de recurso; D troca método sem resolver objetivo; E abandona accountability em vez de melhorar o desenho do indicador.
\end{solucao}

\begin{discursiva}[governança]
O TCE possui PDTI com dezenas de iniciativas, mas sem vínculo explícito a objetivos institucionais, responsáveis, indicadores ou riscos. Em até 30 linhas, descreva como estruturar a governança desse portfólio à luz de COBIT, planejamento estratégico e medição de desempenho.
\end{discursiva}
